% =====================================================================
%  10. RECOMENDACIONES
%  Responsable: INTEGRANTE 5
%  ESTADO: REDACTADO (1 de agosto de 2026) — ~800 palabras
%
%  DIFERENCIA CON LAS CONCLUSIONES: la conclusión mira hacia ATRÁS (qué se
%  encontró); la recomendación mira hacia ADELANTE (qué debe hacerse).
%  Cada recomendación tiene DESTINATARIO explícito: sin destinatario no es
%  una recomendación, es una opinión.
% =====================================================================

\section{Recomendaciones}
\label{sec:recomendaciones}

% ---------------------------------------------------------------------
\subsection{A la empresa operadora y al sector minero}

\begin{enumerate}
  \item \textbf{Establecer la línea base ambiental antes de operar, y hacerla
        pública.} Es la recomendación de mayor rendimiento de todo el informe y
        la de menor costo relativo. Una línea base pública, con puntos de muestreo
        acordados con las comunidades, protege tanto a la población frente a
        impactos no detectados como a la propia empresa frente a atribuciones
        infundadas. Su ausencia en este caso hizo que ninguna de las dos partes
        pudiera demostrar su posición.
  \item \textbf{Someterse a verificación independiente de tercera parte bajo un
        estándar auditable}, y no únicamente a autoevaluación o a evaluaciones de
        encargo. El estándar IRMA ofrece requisitos predefinidos, umbrales
        cuantificados y verificación periódica: los tres elementos que faltaron en
        el proceso analizado.
  \item \textbf{Definir por escrito el criterio de cierre de las no
        conformidades} y someterlo a validación de un tercero. Un plan de acción
        cuyo cumplimiento verifica quien debe cumplirlo no cierra el ciclo de
        mejora continua.
  \item \textbf{Adoptar formalmente los Principios Voluntarios en Seguridad y
        Derechos Humanos} y extender su aplicación a los contratistas de
        seguridad, con registro público de incidentes.
  \item \textbf{Sostener el monitoreo participativo y la transparencia de datos
        durante todo el post-cierre.} El empleo dura lo que dura la mina; el
        pasivo ambiental, no.
\end{enumerate}

% ---------------------------------------------------------------------
\subsection{Al Estado y a los organismos reguladores}

\begin{enumerate}
  \item \textbf{Garantizar la consulta previa, libre e informada antes del
        otorgamiento de la licencia}, no como respuesta al conflicto ya
        desatado. El caso demuestra que una consulta omitida no puede subsanarse
        después: solo puede repararse, que es un remedio inferior y más costoso
        para todas las partes.
  \item \textbf{Exigir línea base ambiental pública como requisito de la
        certificación ambiental}, con puntos de muestreo, parámetros y
        laboratorios acreditados definidos por la autoridad y no por el titular.
  \item \textbf{Dotar a la fiscalización ambiental de capacidad sancionadora
        efectiva y de recursos para ejercerla.} El contraste con el régimen
        peruano muestra que disponer del instrumento no equivale a usarlo con
        oportunidad; la experiencia de Espinar lo confirma.
  \item \textbf{Exigir el cumplimiento del Global Industry Standard on Tailings
        Management} en los depósitos de relaves, incluida la designación del
        Ingeniero de Registro y la constitución del panel independiente de
        revisión técnica.
  \item \textbf{Establecer garantías financieras de cierre suficientes y
        actualizables}, calculadas sobre el horizonte real del pasivo y no sobre
        el de la concesión.
\end{enumerate}

% ---------------------------------------------------------------------
\subsection{A las comunidades y organizaciones sociales}

\begin{enumerate}
  \item \textbf{Participar en los espacios de monitoreo ambiental participativo
        con capacitación técnica propia}, acceso a laboratorios acreditados y
        cadena de custodia documentada. El monitoreo comunitario gana o pierde
        credibilidad por su rigor metodológico, no por su origen: el modelo de
        Espinar es una referencia de escala útil.
  \item \textbf{Documentar y canalizar las quejas por los mecanismos formales},
        conservando registro propio de cada reclamo presentado y de la respuesta
        recibida. Ese registro es, con frecuencia, la única evidencia disponible
        cuando el caso escala a instancias externas.
  \item \textbf{Exigir la publicación de la línea base y de los resultados de
        monitoreo} como condición del diálogo, antes que como reclamo posterior.
\end{enumerate}

% ---------------------------------------------------------------------
\subsection{A la comunidad académica}

\begin{enumerate}
  \item \textbf{Desarrollar estudios epidemiológicos y de biomonitoreo de largo
        plazo con diseño acordado entre las partes}, capaces de distinguir el
        aporte geogénico del antropogénico. El conflicto de evidencia de este caso
        no se resolvió porque ningún estudio individual podía resolverlo con la
        metodología disponible.
  \item \textbf{Investigar comparativamente la eficacia de la auditoría
        voluntaria frente a la fiscalización estatal} en contextos de
        conflictividad socioambiental minera en América Latina, midiendo cambio
        efectivo de conducta y no solo emisión de recomendaciones.
  \item \textbf{Documentar los procesos de post-cierre minero}, que constituyen
        la fase menos estudiada y la de mayor horizonte temporal.
\end{enumerate}
